
\noindent This work is framed as a bounded orbital-dynamics study of a specific short-horizon regime in near-polar LEO in which solar radiation pressure (SRP) can, under restricted passive conditions, produce a net orbit-averaged energy gain. The system considered is intentionally limited to the coupled effects of central gravity, Earth’s oblateness (\(J_2\)), atmospheric drag, solar radiation pressure, and eclipse-modulated shadowing, for a spin-stabilised sail whose orientation is fixed in inertial space and not actively controlled.

\vspace{0.5em}
\par\noindent The primary question is whether eclipse timing can, over a bounded horizon, produce a net positive orbit-averaged energy input by selectively suppressing phases of negative SRP work, and how strongly that effect depends on orbital geometry. The purpose is therefore to isolate and explain a geometric dynamical mechanism, not to demonstrate an operational orbit-raising strategy.

\vspace{0.5em}
\par\noindent The study is framed around two explicit objectives. The first (scientific) objective is to \emph{verify the existence} of the passive-gain region: to map the constrained dynamical regime, explain its underlying eclipse-modulated SRP work-imbalance mechanism, and determine the conditions under which the observed energy gain emerges and subsequently breaks down. The second (methodological) objective is to \emph{validate the open-source propagator} used for the first: to cross-check its force models, component by component, against the industrial flight-dynamics tool STK Astrogator and against the independent open-source library Orekit, and to quantify where the tools agree and diverge. The two objectives are complementary --- the validation establishes that the gain region is reproducible across independent tools and identifies which force-model choices most strongly govern the result.

\vspace{0.5em}
\par\noindent Several exclusions are imposed to maintain a controlled scope. This study does not claim disposal-compliant mission design, communications-feasible spacecraft architecture, collision-risk-assessed operations, maintenance of an exact Sun-synchronous orbit, or long-term sustained orbit raising. It does not attempt subsystem-level mission design, active attitude-control implementation, or guidance and feedback control. The external validation performed here is a component-level cross-check of the propagator's force models against STK Astrogator and Orekit, not a flight-data validation of the mission concept; the results should therefore not be interpreted as a demonstrated mission concept, but as evidence for a bounded dynamical effect within the modelled regime, together with a verified toolchain for studying it. Sun-synchronous-like initial conditions are used only to impose a controlled and interpretable Sun-relative geometry; they are not intended to assert that the vehicle remains on an exactly maintained operational SSO throughout the propagation.
\FloatBarrier